% Options for packages loaded elsewhere
\PassOptionsToPackage{unicode}{hyperref}
\PassOptionsToPackage{hyphens}{url}
\documentclass[
]{article}
\usepackage{xcolor}
\usepackage[margin=1in]{geometry}
\usepackage{amsmath,amssymb}
\setcounter{secnumdepth}{-\maxdimen} % remove section numbering
\usepackage{iftex}
\ifPDFTeX
  \usepackage[T1]{fontenc}
  \usepackage[utf8]{inputenc}
  \usepackage{textcomp} % provide euro and other symbols
\else % if luatex or xetex
  \usepackage{unicode-math} % this also loads fontspec
  \defaultfontfeatures{Scale=MatchLowercase}
  \defaultfontfeatures[\rmfamily]{Ligatures=TeX,Scale=1}
\fi
\usepackage{lmodern}
\ifPDFTeX\else
  % xetex/luatex font selection
\fi
% Use upquote if available, for straight quotes in verbatim environments
\IfFileExists{upquote.sty}{\usepackage{upquote}}{}
\IfFileExists{microtype.sty}{% use microtype if available
  \usepackage[]{microtype}
  \UseMicrotypeSet[protrusion]{basicmath} % disable protrusion for tt fonts
}{}
\makeatletter
\@ifundefined{KOMAClassName}{% if non-KOMA class
  \IfFileExists{parskip.sty}{%
    \usepackage{parskip}
  }{% else
    \setlength{\parindent}{0pt}
    \setlength{\parskip}{6pt plus 2pt minus 1pt}}
}{% if KOMA class
  \KOMAoptions{parskip=half}}
\makeatother
\usepackage{color}
\usepackage{fancyvrb}
\newcommand{\VerbBar}{|}
\newcommand{\VERB}{\Verb[commandchars=\\\{\}]}
\DefineVerbatimEnvironment{Highlighting}{Verbatim}{commandchars=\\\{\}}
% Add ',fontsize=\small' for more characters per line
\usepackage{framed}
\definecolor{shadecolor}{RGB}{248,248,248}
\newenvironment{Shaded}{\begin{snugshade}}{\end{snugshade}}
\newcommand{\AlertTok}[1]{\textcolor[rgb]{0.94,0.16,0.16}{#1}}
\newcommand{\AnnotationTok}[1]{\textcolor[rgb]{0.56,0.35,0.01}{\textbf{\textit{#1}}}}
\newcommand{\AttributeTok}[1]{\textcolor[rgb]{0.13,0.29,0.53}{#1}}
\newcommand{\BaseNTok}[1]{\textcolor[rgb]{0.00,0.00,0.81}{#1}}
\newcommand{\BuiltInTok}[1]{#1}
\newcommand{\CharTok}[1]{\textcolor[rgb]{0.31,0.60,0.02}{#1}}
\newcommand{\CommentTok}[1]{\textcolor[rgb]{0.56,0.35,0.01}{\textit{#1}}}
\newcommand{\CommentVarTok}[1]{\textcolor[rgb]{0.56,0.35,0.01}{\textbf{\textit{#1}}}}
\newcommand{\ConstantTok}[1]{\textcolor[rgb]{0.56,0.35,0.01}{#1}}
\newcommand{\ControlFlowTok}[1]{\textcolor[rgb]{0.13,0.29,0.53}{\textbf{#1}}}
\newcommand{\DataTypeTok}[1]{\textcolor[rgb]{0.13,0.29,0.53}{#1}}
\newcommand{\DecValTok}[1]{\textcolor[rgb]{0.00,0.00,0.81}{#1}}
\newcommand{\DocumentationTok}[1]{\textcolor[rgb]{0.56,0.35,0.01}{\textbf{\textit{#1}}}}
\newcommand{\ErrorTok}[1]{\textcolor[rgb]{0.64,0.00,0.00}{\textbf{#1}}}
\newcommand{\ExtensionTok}[1]{#1}
\newcommand{\FloatTok}[1]{\textcolor[rgb]{0.00,0.00,0.81}{#1}}
\newcommand{\FunctionTok}[1]{\textcolor[rgb]{0.13,0.29,0.53}{\textbf{#1}}}
\newcommand{\ImportTok}[1]{#1}
\newcommand{\InformationTok}[1]{\textcolor[rgb]{0.56,0.35,0.01}{\textbf{\textit{#1}}}}
\newcommand{\KeywordTok}[1]{\textcolor[rgb]{0.13,0.29,0.53}{\textbf{#1}}}
\newcommand{\NormalTok}[1]{#1}
\newcommand{\OperatorTok}[1]{\textcolor[rgb]{0.81,0.36,0.00}{\textbf{#1}}}
\newcommand{\OtherTok}[1]{\textcolor[rgb]{0.56,0.35,0.01}{#1}}
\newcommand{\PreprocessorTok}[1]{\textcolor[rgb]{0.56,0.35,0.01}{\textit{#1}}}
\newcommand{\RegionMarkerTok}[1]{#1}
\newcommand{\SpecialCharTok}[1]{\textcolor[rgb]{0.81,0.36,0.00}{\textbf{#1}}}
\newcommand{\SpecialStringTok}[1]{\textcolor[rgb]{0.31,0.60,0.02}{#1}}
\newcommand{\StringTok}[1]{\textcolor[rgb]{0.31,0.60,0.02}{#1}}
\newcommand{\VariableTok}[1]{\textcolor[rgb]{0.00,0.00,0.00}{#1}}
\newcommand{\VerbatimStringTok}[1]{\textcolor[rgb]{0.31,0.60,0.02}{#1}}
\newcommand{\WarningTok}[1]{\textcolor[rgb]{0.56,0.35,0.01}{\textbf{\textit{#1}}}}
\usepackage{graphicx}
\makeatletter
\newsavebox\pandoc@box
\newcommand*\pandocbounded[1]{% scales image to fit in text height/width
  \sbox\pandoc@box{#1}%
  \Gscale@div\@tempa{\textheight}{\dimexpr\ht\pandoc@box+\dp\pandoc@box\relax}%
  \Gscale@div\@tempb{\linewidth}{\wd\pandoc@box}%
  \ifdim\@tempb\p@<\@tempa\p@\let\@tempa\@tempb\fi% select the smaller of both
  \ifdim\@tempa\p@<\p@\scalebox{\@tempa}{\usebox\pandoc@box}%
  \else\usebox{\pandoc@box}%
  \fi%
}
% Set default figure placement to htbp
\def\fps@figure{htbp}
\makeatother
\setlength{\emergencystretch}{3em} % prevent overfull lines
\providecommand{\tightlist}{%
  \setlength{\itemsep}{0pt}\setlength{\parskip}{0pt}}
\usepackage{bookmark}
\IfFileExists{xurl.sty}{\usepackage{xurl}}{} % add URL line breaks if available
\urlstyle{same}
\hypersetup{
  pdftitle={Propensity scores and IPTW},
  hidelinks,
  pdfcreator={LaTeX via pandoc}}

\title{Propensity scores and IPTW}
\author{}
\date{\vspace{-2.5em}}

\begin{document}
\maketitle

\textbf{Inference lab} using the five-step template: Hypothesis →
Visualise → Assumptions → Conduct → Conclude.

\subsection{Learning objectives}\label{learning-objectives}

\begin{itemize}
\tightlist
\item
  Estimate a propensity score and match nearest-neighbour with
  \texttt{MatchIt}.
\item
  Construct inverse-probability-of-treatment weights (IPTW) and fit a
  weighted outcome model.
\item
  Diagnose covariate balance with \texttt{cobalt::love.plot}.
\end{itemize}

\subsection{Prerequisites}\label{prerequisites}

DAGs and logistic regression.

\subsection{Background}\label{background}

The propensity score is the probability of being treated given
covariates. Under the assumptions of conditional exchangeability and
positivity, conditioning on the propensity score yields an unbiased
estimate of the average treatment effect. Matching and
inverse-probability-of-treatment weighting are two ways to do that
conditioning.

Matching pairs each treated unit with one (or more) untreated units with
similar propensity scores, then analyses the matched sample. IPTW gives
each unit a weight of 1/e(X) if treated and 1/(1 − e(X)) if untreated,
creating a pseudo-population where treatment is independent of X. Both
estimators require the propensity model to be correctly specified;
neither handles unmeasured confounding.

Balance diagnostics assess whether the matched or weighted sample has
covariate distributions that are similar across treatment groups. The
standardised mean difference (SMD) is the workhorse metric; anything
below 0.1 is usually considered balanced.

Large weights (extreme propensities near 0 or 1) are a sign of poor
positivity. Stabilised weights and truncation can help, but the
underlying problem --- a region of covariate space where one treatment
is never seen --- cannot be fixed statistically.

\subsection{Setup}\label{setup}

\begin{Shaded}
\begin{Highlighting}[]
\FunctionTok{library}\NormalTok{(tidyverse)}
\FunctionTok{library}\NormalTok{(MatchIt)}
\FunctionTok{library}\NormalTok{(cobalt)}
\FunctionTok{set.seed}\NormalTok{(}\DecValTok{42}\NormalTok{)}
\FunctionTok{theme\_set}\NormalTok{(}\FunctionTok{theme\_minimal}\NormalTok{(}\AttributeTok{base\_size =} \DecValTok{12}\NormalTok{))}
\end{Highlighting}
\end{Shaded}

\subsection{1. Hypothesis}\label{hypothesis}

Treatment reduces the outcome by a known amount. We will estimate the
effect with naive regression, matching, and IPTW, and compare with the
simulated truth.

\subsection{2. Visualise}\label{visualise}

\begin{Shaded}
\begin{Highlighting}[]
\NormalTok{dat }\OtherTok{\textless{}{-}} \FunctionTok{tibble}\NormalTok{(}
  \AttributeTok{age  =} \FunctionTok{rnorm}\NormalTok{(n, }\DecValTok{60}\NormalTok{, }\DecValTok{10}\NormalTok{),}
  \AttributeTok{sev  =} \FunctionTok{rnorm}\NormalTok{(n, }\DecValTok{0}\NormalTok{,   }\DecValTok{1}\NormalTok{),}
  \AttributeTok{sex  =} \FunctionTok{rbinom}\NormalTok{(n, }\DecValTok{1}\NormalTok{, }\FloatTok{0.5}\NormalTok{)}
\NormalTok{) }\SpecialCharTok{|\textgreater{}}
  \FunctionTok{mutate}\NormalTok{(}\AttributeTok{ps =} \FunctionTok{plogis}\NormalTok{(}\SpecialCharTok{{-}}\DecValTok{1} \SpecialCharTok{+} \FloatTok{0.04} \SpecialCharTok{*}\NormalTok{ (age }\SpecialCharTok{{-}} \DecValTok{60}\NormalTok{) }\SpecialCharTok{+} \FloatTok{0.8} \SpecialCharTok{*}\NormalTok{ sev }\SpecialCharTok{+} \FloatTok{0.3} \SpecialCharTok{*}\NormalTok{ sex),}
         \AttributeTok{trt =} \FunctionTok{rbinom}\NormalTok{(n, }\DecValTok{1}\NormalTok{, ps),}
         \AttributeTok{y   =} \DecValTok{2} \SpecialCharTok{{-}} \FloatTok{1.5} \SpecialCharTok{*}\NormalTok{ trt }\SpecialCharTok{+} \FloatTok{0.05} \SpecialCharTok{*}\NormalTok{ (age }\SpecialCharTok{{-}} \DecValTok{60}\NormalTok{) }\SpecialCharTok{+}
                 \FloatTok{0.8} \SpecialCharTok{*}\NormalTok{ sev }\SpecialCharTok{+} \FloatTok{0.3} \SpecialCharTok{*}\NormalTok{ sex }\SpecialCharTok{+} \FunctionTok{rnorm}\NormalTok{(n, }\DecValTok{0}\NormalTok{, }\DecValTok{1}\NormalTok{))}

\FunctionTok{ggplot}\NormalTok{(dat, }\FunctionTok{aes}\NormalTok{(ps, }\AttributeTok{fill =} \FunctionTok{factor}\NormalTok{(trt))) }\SpecialCharTok{+}
  \FunctionTok{geom\_density}\NormalTok{(}\AttributeTok{alpha =} \FloatTok{0.5}\NormalTok{) }\SpecialCharTok{+}
  \FunctionTok{labs}\NormalTok{(}\AttributeTok{x =} \StringTok{"Propensity score"}\NormalTok{, }\AttributeTok{y =} \StringTok{"Density"}\NormalTok{, }\AttributeTok{fill =} \StringTok{"Treated?"}\NormalTok{)}
\end{Highlighting}
\end{Shaded}

\subsection{3. Assumptions}\label{assumptions}

No unmeasured confounding (conditional exchangeability); positivity
(everyone has a non-zero chance of each treatment); correct
specification of the propensity model.

\subsection{4. Conduct}\label{conduct}

\begin{Shaded}
\begin{Highlighting}[]
\NormalTok{fit\_naive }\OtherTok{\textless{}{-}} \FunctionTok{lm}\NormalTok{(y }\SpecialCharTok{\textasciitilde{}}\NormalTok{ trt }\SpecialCharTok{+}\NormalTok{ age }\SpecialCharTok{+}\NormalTok{ sev }\SpecialCharTok{+}\NormalTok{ sex, }\AttributeTok{data =}\NormalTok{ dat)}
\FunctionTok{coef}\NormalTok{(fit\_naive)[}\StringTok{"trt"}\NormalTok{]}
\end{Highlighting}
\end{Shaded}

\begin{Shaded}
\begin{Highlighting}[]
\NormalTok{m }\OtherTok{\textless{}{-}} \FunctionTok{matchit}\NormalTok{(trt }\SpecialCharTok{\textasciitilde{}}\NormalTok{ age }\SpecialCharTok{+}\NormalTok{ sev }\SpecialCharTok{+}\NormalTok{ sex, }\AttributeTok{data =}\NormalTok{ dat,}
             \AttributeTok{method =} \StringTok{"nearest"}\NormalTok{, }\AttributeTok{ratio =} \DecValTok{1}\NormalTok{)}
\NormalTok{m}
\NormalTok{matched }\OtherTok{\textless{}{-}} \FunctionTok{match.data}\NormalTok{(m)}

\NormalTok{fit\_match }\OtherTok{\textless{}{-}} \FunctionTok{lm}\NormalTok{(y }\SpecialCharTok{\textasciitilde{}}\NormalTok{ trt, }\AttributeTok{data =}\NormalTok{ matched,}
                \AttributeTok{weights =}\NormalTok{ matched}\SpecialCharTok{$}\NormalTok{weights)}
\FunctionTok{coef}\NormalTok{(fit\_match)[}\StringTok{"trt"}\NormalTok{]}
\end{Highlighting}
\end{Shaded}

\begin{Shaded}
\begin{Highlighting}[]
\NormalTok{dat }\OtherTok{\textless{}{-}}\NormalTok{ dat }\SpecialCharTok{|\textgreater{}}
  \FunctionTok{mutate}\NormalTok{(}\AttributeTok{w =} \FunctionTok{if\_else}\NormalTok{(trt }\SpecialCharTok{==} \DecValTok{1}\NormalTok{, }\DecValTok{1} \SpecialCharTok{/}\NormalTok{ ps, }\DecValTok{1} \SpecialCharTok{/}\NormalTok{ (}\DecValTok{1} \SpecialCharTok{{-}}\NormalTok{ ps)))}
\NormalTok{fit\_iptw }\OtherTok{\textless{}{-}} \FunctionTok{lm}\NormalTok{(y }\SpecialCharTok{\textasciitilde{}}\NormalTok{ trt, }\AttributeTok{data =}\NormalTok{ dat, }\AttributeTok{weights =}\NormalTok{ w)}
\FunctionTok{coef}\NormalTok{(fit\_iptw)[}\StringTok{"trt"}\NormalTok{]}
\end{Highlighting}
\end{Shaded}

\begin{Shaded}
\begin{Highlighting}[]
\FunctionTok{love.plot}\NormalTok{(m, }\AttributeTok{thresholds =} \FunctionTok{c}\NormalTok{(}\AttributeTok{m =} \FloatTok{0.1}\NormalTok{), }\AttributeTok{abs =} \ConstantTok{TRUE}\NormalTok{)}
\end{Highlighting}
\end{Shaded}

\subsection{5. Concluding statement}\label{concluding-statement}

\begin{quote}
The simulated treatment effect was −1.5. The naive regression gave
\texttt{round(coef(fit\_naive){[}"trt"{]},\ 2)}; nearest-neighbour
matching gave \texttt{round(coef(fit\_match){[}"trt"{]},\ 2)}; IPTW gave
\texttt{round(coef(fit\_iptw){[}"trt"{]},\ 2)}. Balance plots showed
that matching reduced all SMDs below 0.1.
\end{quote}

\subsection{Common pitfalls}\label{common-pitfalls}

\begin{itemize}
\tightlist
\item
  Matching with replacement without adjusting inference for the repeated
  use of controls.
\item
  Reporting only the \emph{p}-value for the treatment effect in the
  weighted model without a sandwich or bootstrap standard error.
\item
  Ignoring extreme weights.
\item
  Picking the matching method that gives the effect you wanted.
\end{itemize}

\subsection{Further reading}\label{further-reading}

\begin{itemize}
\tightlist
\item
  Austin PC (2011), \emph{An introduction to propensity score methods
  for reducing the effects of confounding in observational studies}.
\item
  Stuart EA (2010), \emph{Matching methods for causal inference}.
\item
  Hernán MA, Robins JM. \emph{Causal Inference: What If}, ch.~12.
\end{itemize}

\subsection{Session info}\label{session-info}

\begin{Shaded}
\begin{Highlighting}[]

\end{Highlighting}
\end{Shaded}

\subsection{See also --- chapter index}\label{see-also-chapter-index}

\begin{itemize}
\tightlist
\item
  \href{../index.Rmd}{Methods in this chapter}
\item
  \href{../../../ch00-find-your-method.Rmd}{Find your method (Chapter
  0)}
\end{itemize}

\end{document}
